\vekalet{\textbf{Two base and aligned continuations of one prompt, chosen by a rule written before the choice.} For each of two models we computed $\Delta M_1$ per prompt on the raw-continuation debate prompts (all sixteen prefix conditions and twelve draws), took the prompt whose $\Delta M_1$ lies closest to the model's median over its $34$ prompts, and printed the first prefix condition without a second person of its own (ARO-x1 in Table~\ref{tab:onek}), first draw. The rule skips a prompt whose text contains listed offensive terms or would trip our anonymity or vocabulary checks, and says so; each text is cut at its first $120$ tokens (the counter's own tokens) and printed verbatim, line breaks shown as $\hookleftarrow$. Hashes are SHA-256 prefixes of the printed text and of the generation file it comes from. The panel assigns each thesis both stances; a continuation defends the stance it was given, and the example is the rule-selected prompt, not a chosen one (the selection rule is in the released repository, committed before the selection).}

\paragraph*{Mistral-7B-v0.3.}
\vekalet{Model median $\Delta M_1$ $-8.95$ per thousand; prompt 2, $\Delta M_1$ $-9.17$; no candidate was skipped. Generation files: base \texttt{f5291d44eb5e37c9}, aligned \texttt{e4b5e106f119b9fc}. The two printed windows carry $0$ and $0$ first-person and $0$ and $0$ second-person forms (base, aligned): the rule picks the prompt at the model's median, not a window that shows the fall; Figure~\ref{fig:ornek} prints the Gemma-3-27B pair, where the base window has $2$ first-person forms and the aligned window none.}
\begin{quote}\small\raggedright
\textbf{Prompt} (\texttt{43cb3a569a6be153}): Pineapple belongs on pizza is right. Sweet and salty combinations occur in several cuisines. Preferences differ between people. The topping is sold in many places. That is the position. bullfighting as a male dominated sport reinforces harmful gender stereotypes is right. Here is why:\par\smallskip
\textbf{Base} (\texttt{903fa77b9a5ffa14}): Gender stereotypes are an aspect of a patriarchal system that causes harm to people, especially women and the LGBTQIA community. $\hookleftarrow$ But that's not the only way bullfighting reinforces harmful gender stereotypes. It also promotes the idea that men are strong and brave while women are fragile and weak. $\hookleftarrow$ A few decades ago, bullfighting was a prominent event. But with the rise of other forms of entertainment, bullfighting has now become a thing of the past. But some people still enjoy bullfighting. In fact, there are even bullfighting shows that are held in some parts of the world. Some people argue that bullfighting is a part of the country\par\smallskip
\textbf{Aligned} (\texttt{03e7dcb9c2ee4fcf}):  $\hookleftarrow$   1. Domination: Bullfighting involves one party (the man) dominating another (the bull), which reinforces traditional masculine roles of dominance and power. $\hookleftarrow$   2. Violence: The violence inherent in bullfighting is associated with masculinity and can perpetuate harmful ideas about the use of force to assert dominance. $\hookleftarrow$   3. Spectacle: The spectacle of bullfighting is often associated with macho behavior and can reinforce the idea that violence is entertainment, especially for men. $\hookleftarrow$   4. Limited Participation: The sport is largely restricted to men, limiting opportunities for women and reinforcing gender roles. $\hookleftarrow$   5. Cultural Perception: The continued acceptance and glorification of bullfighting perpet
\end{quote}

\paragraph*{Gemma-3-27B.}
\vekalet{Model median $\Delta M_1$ $-2.15$ per thousand; prompt 7, $\Delta M_1$ $-1.80$; no candidate was skipped. Generation files: base \texttt{0a74fa0aec1f728b}, aligned \texttt{0a62776eddf16eed}.}
\begin{quote}\small\raggedright
\textbf{Prompt} (\texttt{fd18dd6401ea0336}): Pineapple belongs on pizza is right. Sweet and salty combinations occur in several cuisines. Preferences differ between people. The topping is sold in many places. That is the position. feminism aims to correct the systemic biases against women and other minorities that are born out of historical discrimination is wrong. Here is why:\par\smallskip
\textbf{Base} (\texttt{7cdd4470fae4dd3d}):  I did a quick search of the word feminism in the Oxford dictionary, and these were the definitions I found. 1. The advocacy of women’s rights on the ground of the equality of the sexes. 2. The theory of the political, economic, and social equality of the sexes. 3. Organized activity on behalf of women’s rights and interests. All of these definitions are correct, and all of them are what is meant when someone says they are a feminist. Here are some of the other things that feminism is not: feminism is not about hating men. Feminism is not about being better than men\par\smallskip
\textbf{Aligned} (\texttt{ac3e8d78edbc6a86}):  $\hookleftarrow$ 1.  **Feminism's Scope Has Expanded Beyond Addressing Discrimination:** While initially focused on equal rights and opportunities, modern feminism often extends into areas beyond addressing demonstrable discrimination. It frequently includes advocating for specific social outcomes, challenging traditional gender roles, and analyzing societal structures through a particular ideological lens. This expanded scope goes beyond simply correcting biases and enters the realm of social engineering. $\hookleftarrow$ 2.  **Defining "Systemic Bias" is Subjective:** The concept of "systemic bias" is often difficult to objectively define and measure. What one person perceives as systemic bias, another may see as the natural outcome of individual
\end{quote}
